\section{函数}
从静止开始的自由落体，其下落距离 \(s\) 与时间 \(t\) 的关系是 \(s=\frac12 gt^2\)；质量为 \(m\) 的质点，其动能 \(E\) 与速率 \(v\) 的关系是 \(E=\frac12 mv^2\)；直角三角形的一个锐角固定为 \(\alpha\)，若与它相邻的直角边长为 \(x\)，则面积为 \(S=\frac12(\tan\alpha)x^2\)。三个问题中的量各有不同，所取的函数却都有 \(y=kx^2\) 的形式。

哪些量作为自变量，取决于所问的问题。例如，不再固定质点的质量时，动能便是质量和速率这两个量的函数：给定一对 \((m,v)\)，就确定一个数 \(E\)。反过来，函数的值也可以是向量。设一个质点沿抛物线 \(y=x^2\) 运动，在时刻 \(t\) 的横坐标为 \(t\)。若问它有多高，答案是 \(y(t)=t^2\)；若问它在哪里，答案是位置向量 \(\mathbf r(t)=(t,t^2)\)。后者称为向量值函数：两个坐标合起来才是这个时刻的一个函数值。下面先讨论一个实变量的实值函数。

\subsection{实值函数}
对于定义在实数集合\(\mathbb{R}\)的子集上且取值为实数的函数，严格定义如下：

\begin{definition}
设\(A\)是\(\mathbb{R}\)的子集，若对\(A\)中的每一个数\(x\)，都有唯一确定的\(y \in \mathbb{R}\)与之对应（即函数的单值性），将\(y\)记为\(f(x)\)，那么，就称\(f\)是\(A\)上的一个{\heiti{}实值函数}。集合\(A\)称为\(f\)的定义域，而数\(f(x)\)称为\(f\)的值。\(f\)的一切值的集合叫做\(f\)的值域，通常记作\(f(A)\)，即\(f(A) = \left\{ y \mid y = f(x) , x \in A \right\}\)。习惯上，称上述的\(x\)为自变量，\(y\)为因变量。
\end{definition}

根据上述定义，容易知道函数有3个基本要素，即定义域、值域和（单值的）对应关系。一个函数也可以看成是一个将集合\(A \subset \mathbb{R}\)映入\(\mathbb{R}\)的一个映射：
\begin{equation}
f : A \longrightarrow \mathbb{R} \text{或} f : x \longmapsto y = f(x)
\end{equation}

例如，\(f(x)=x^2\) 在 \([-1,1]\) 上的值域是 \([0,1]\)。若把定义域改成 \([0,1]\)，值域不变，得到的却是另一个函数：原来 \(x=-\frac12\) 和 \(x=\frac12\) 都对应于 \(\frac14\)，现在只有后者。可见，知道一个函数怎样取值，还须知道它在哪里取值。

设函数的定义域为\(A\)，一般而言，其值域\(f(A)\)的性态极其复杂，但有三种情形是值得注意的。

\newpage
1°函数的值域\(f(A)\)是一个有界数集。若存在\(M\geq 0\)，使得\(f(A)\subseteq[-M,M]\)，即对一切\(x\in A\)都有\(\lvert f(x)\rvert\leq M\)，则称\(f\)为\(A\)上的有界函数。

2°定义域\(A\)与值域\(f(A)\)同序（或者反序）。即\(A\)中任意两个数\(x_{1} , x_{2}\)的大小顺序，均与它们对应的值域\(f(A)\)中的两个数\(y_{1} = f(x_{1}) , y_{2} = f(x_{2})\)的大小顺序相同（或者相反），此时称函数\(f(x)\)是单调的。确切地说，\(f(x)\)称为\(A\)上的\\
\hspace*{\parindent} 单调增函数，对任意的\(x_{1} , x_{2} \in A\)，如果\(x_{1} < x_{2}\)，有\(f(x_{1}) \le f(x_{2})\)；\\
\hspace*{\parindent} 单调减函数，对任意的\(x_{1} , x_{2} \in A\)，如果\(x_{1} < x_{2}\)，有\(f(x_{1}) \ge f(x_{2})\)；

单调增和单调减函数，统称为单调函数。若上面的不等号为严格不等号，则称\(f(x)\)为严格单调增（减）函数。

3°定义域\(A\)与值域\(f(A)\)一一对应。即对每一个\(y \in f(A)\)都有唯一确定的\(x \in A\)使得\(f(x) = y\)。此时，自然地可以导出一个由\(f(A)\)到\(A\)的映射。这个映射称为\(f(x)\)的反函数（或逆映射），记为\(f^{- 1}\)，即\(x = f^{- 1}(y)\)。它的定义域为\(f(A)\)，值域为\(A\)。

显然，若\(f\)是\(A\)上严格单调的函数，则\(f\)必然是\(A\)到\(f(A)\)的一个一一映射，因此也就有反函数。

回看上面的平方函数。在 \([-1,1]\) 上，由平方值不能唯一恢复原数；限制在 \([0,1]\) 上，它严格递增，反函数为 \(y=\sqrt{x}\)。若保留的恰好是左半段 \([-1,0]\)，反函数又是什么？

\subsection{函数的复合}
不同的函数在它们公共有定义的区域内是可以进行四则运算的。这里主要介绍函数的另一种运算——函数的复合运算。

设有函数\(y = f(u)\)，定义域为\(B\)，值域为\(C\)，\(u = g(x)\)，定义域为\(A\)。若\(g(A)\subseteq B\)，那么由
\begin{equation}
f \circ g(x) = f(g(x))
\end{equation}
定义了一个新函数\(f \circ g\)，称之为\(f\)与\(g\)的复合函数，它的定义域是\(A\)，通常记为\(y = f(g(x))\)，而称\(u\)为中间变量。
\begin{center}
\begin{minipage}{0.92\linewidth}
\setlength{\parskip}{0pt}
\centering
\includegraphics[width=\linewidth]{function-composition}
\setlength{\abovecaptionskip}{3pt}
\captionof{figure}{函数的复合}\label{fig:function-composition}
\end{minipage}
\end{center}
从映射的角度看，图\ref{fig:function-composition}中的复合函数就是一个从\(A\)到\(B\)再到\(C\)的映射。

取 \(f(u)=\sqrt u\)，\(g(x)=x^2\)。从任意实数出发，先平方再开平方，结果是 \(|x|\)；先开平方再平方，只能从非负数出发，结果是 \(x\)。也就是说，
\[
 (f\circ g)(x)=|x| \quad (x\in\mathbb R),\qquad
 (g\circ f)(x)=x \quad (x\geq 0).
\]
这里复合的顺序既影响对应法则，也影响可以从哪些数出发。

注意到，\(f \circ g\)与\(g \circ f\)一般并不相同，但是容易看出，复合运算满足结合律
\begin{equation}
\left( f \circ g \right) \circ h = f \circ \left( g \circ h \right)
\end{equation}
因此可以用\(f \circ g \circ h\)表示，或者记为\(f(g(h(x)))\)，这是3个函数的复合。类似地可以考虑任意有限多个函数的复合。特别地，记\(n\)个\(f\)复合为\(f \circ f \circ \cdots \circ f = f^{n}(x)\)。

{\kaiti{}思考\quad 分析上述复合函数\(f \circ g\)的值域如何确定。}

{\kaiti{}\hspace*{\parindent}\(f\)和\(g\)满足什么条件时，\(f \circ g\)与\(g \circ f\)是相同的。}


\subsection{初等函数}
初等函数是一类简单的函数，有简单而良好的性态。其中最基本的函数是多项式函数、幂函数、指数函数、对数函数、三角函数与反三角函数，称它们为{\heiti{}基本初等函数}。由基本初等函数经过有限次加减乘除和复合运算得出的函数称为{\heiti{}初等函数}。这里简单介绍几种初等函数，初等函数的具体性质在后续内容中进行分析。

设\(n\)是非负整数，形如
\begin{equation}
f(x) = a_{n} x^{n} + \cdots + a_{1} x + a_{0}
\end{equation}
的函数称为多项式函数。两个多项式函数\(f(x)\)、\(g(x)\)的商称为有理函数，它的定义域是不包括\(g(x) = 0\)的所有实数。

形如\(f(x) = x^{\alpha}\)的函数称为幂函数，其中\(\alpha\)可以是任意实数。试分析不同\(\alpha\)取值下的\(f(x)\)的性质。

所谓双曲函数是指如下几个函数，它们的定义域是\(\mathbb{R}\)。
\begin{equation}
\sinh x = \frac{e^{x} - e^{- x}}{2} , \cosh x = \frac{e^{x} + e^{- x}}{2} ,
\end{equation}
\begin{equation}
\tanh x = \frac{e^{x} - e^{- x}}{e^{x} + e^{- x}} , \coth x = \frac{e^{x} + e^{- x}}{e^{x} - e^{- x}} ( x \ne 0 )
\end{equation}
由定义可见，\(\cosh t+\sinh t=e^t\)，\(\cosh t-\sinh t=e^{-t}\)。两式相乘，得到 \(\cosh^2 t-\sinh^2 t=1\)。

求双曲正弦函数\(y = \sinh x\)的反函数\(y = \sinh^{- 1} x\).


\subsection{函数的其他表达方式}
通常把函数\(y = f(x)\)的表达方式称为显式表示或者显表示。

\paragraph{1°分段函数} 所谓分段函数，简单地讲，就是对于自变量不同的取值范围，相应的函数表达式（即对应关系）不同。

绝对值函数在数轴上给出一个点到原点的距离。同一条取值规则可以写成两段：
\[
 |x|=\begin{cases}-x,&x<0,\\ x,&x\geq 0.\end{cases}
\]

\paragraph{2°函数的隐表示} 有时自变量与函数值之间的对应关系由一个方程给出。例如，单位圆的方程为 \(x^2+y^2=1\)。给定 \(-1<x<1\)，圆上有两点具有这个横坐标，故整个圆不能写成一个实值函数 \(y=f(x)\) 的图像。若指定取上半圆，方程就唯一确定了 \(y\)；把它解出，即得显式表示 \(y=\sqrt{1-x^2}\)，其中 \(-1\leq x\leq 1\)。取下半圆则得另一个函数。

\paragraph{3°函数的参数方程表示} 所谓参数方程表示，是指曲线上点的两个坐标\(x\)与\(y\)之间的关系是通过一个参数相联系的，即\(x\)和\(y\)可以表示成下列形式
\begin{equation}
\left\{ \begin{matrix}x = x(t) , \\ y = y(t) ,\end{matrix} \quad t \in \left[ \alpha , \beta \right] \right.
\end{equation}

仍取单位圆。让点从 \((1,0)\) 出发，逆时针转过 \(t\) 弧度，则它的位置由一个向量值函数给出：
\[
 \mathbf r(t)=(\cos t,\sin t),\qquad 0\leq t\leq 2\pi .
\]
每个 \(t\) 都确定圆上的一个点，无须把圆分成上下两半。消去 \(t\) 得到 \(x^2+y^2=1\)，圆的形状保留了，点从哪里出发、向哪边走却看不出来。若改取 \(\mathbf r(t)=(\cos t,-\sin t)\)，得到的还是同一个圆，运动方向恰好相反。

回看前面的恒等式 \(\cosh^2 t-\sinh^2 t=1\)。正如 \((\cos t,\sin t)\) 描出单位圆，当 \(t\) 取遍实数时，\((\cosh t,\sinh t)\) 描出双曲线 \(x^2-y^2=1\) 的右支。为什么只得到右支？

\newpage
\subsection{通常研究的函数的性质}
\paragraph{1°连续性} 第一章已经给出了函数极限的定义。极限考察自变量接近某点时，函数值是否接近一个确定的数；连续性进一步要求这个数就是函数在该点的值。观察下面三个函数在\(x=0\)附近的图像：
\begin{center}
\begin{minipage}{0.92\linewidth}
\setlength{\parskip}{0pt}
\centering
\includegraphics[width=\linewidth]{function-continuity}
\setlength{\abovecaptionskip}{3pt}
\captionof{figure}{一点处的间断与连续}\label{fig:function-continuity}
\end{minipage}
\end{center}
图\ref{fig:function-continuity}中的区别来自函数在该点附近的极限与该点函数值的关系。函数在\(x_0\)处连续，应当满足\(\lim_{x\to x_0}f(x)=f(x_0)\)。下一节将进一步研究这一性质及其对函数取值的影响。

\paragraph{2°单调性} 单调函数的定义前面已经给出，但是函数并不总是单调函数，因此将函数的定义域分为数个子区间，使得函数在该区间内成为单调函数是便于分析函数的性态的。仿照单调函数的定义，可以类似地定义函数在某区间上是（严格）单调增（减）的，这个区间称为函数的单调增（减）区间。

\paragraph{3°对称性} 函数图像的对称性可以用函数值之间的关系来描述。设函数\(f\)的定义域\(A\)关于\(a\)对称，即\(a+x\in A\)当且仅当\(a-x\in A\)。以下等式均要求对一切满足\(a+x\in A\)的\(x\)成立。

函数\(y=f(x)\)的图像关于直线\(x=a\)对称，当且仅当\(f(a+x)=f(a-x)\)。

函数\(y=f(x)\)的图像关于点\((a,b)\)对称，当且仅当\(f(a+x)+f(a-x)=2b\)。此时，横坐标为\(a+x\)和\(a-x\)的两点以\((a,b)\)为中点。

{\kaiti\textbf{（奇偶性）}}特别地，图像关于原点对称的函数称为{\heiti{}奇函数}，此时\(f(-x)=-f(x)\)；图像关于\(y\)轴对称的函数称为{\heiti{}偶函数}，此时\(f(-x)=f(x)\)。

\paragraph{4°周期性} 设函数\(f(x)\)的定义域是\(A\)，那么称这个函数是周期函数，如果存在一个非零常数\(T\)，使得对于任意的\(x \in A\)，都有\(x + T \in A\)且\(f(x + T) = f(x)\)。显然满足条件的\(T\)有无数个，因此规定：若其中存在大于\(0\)的最小的\(T\)，则称之为这个周期函数的最小正周期。
